\documentclass[11pt]{article}

\usepackage[margin=1in]{geometry}
\usepackage{amsmath}
\usepackage{amssymb}
\usepackage{booktabs}
\usepackage[colorlinks=true,linkcolor=blue,urlcolor=blue]{hyperref}

\setlength{\parindent}{0pt}
\setlength{\parskip}{0.55em}

\newcommand{\pts}[1]{\hspace{0.35em}\textnormal{\textbf{[#1 points]}}}
\newcommand{\St}{S_T}
\newcommand{\Fzero}{F_0}
\newcommand{\Fone}{F_1}
\newcommand{\Ftwo}{F_2}
\newcommand{\K}{K}

\newenvironment{solution}{\par\noindent\textbf{Solution.}\ }{\par}

\title{S181M116 Derivatives and Alternative Investments\\
Assignment 1: Solution Key}
\author{Based on Lecture 1 and Lecture 2, Sections 1--3}
\date{}

\begin{document}
\maketitle

\textbf{Total: 20 points.} This solution key corresponds to the Moodle
assignment on derivative basics, futures markets, hedging, basis risk, and
forward/futures pricing. It excludes the later Lecture 2 material on commodity
behavior, stock-index futures and beta hedging, and stack-and-roll hedging.

\section*{Question 1: Concepts, Market Structure, and Risk Controls \pts{3}}

\begin{solution}
\begin{enumerate}
    \item
    \begin{enumerate}
        \item This is a derivative. The underlying variable is the terminal asset
        price $\St$.
        \item This is not a derivative. It is the underlying security itself.
        \item This is a derivative. The underlying variable is the EUR/USD
        exchange rate.
        \item This is a derivative. The underlying variable is the July
        temperature index or average temperature measure specified in the
        contract.
    \end{enumerate}

    \item The natural starting point is an OTC forward contract. The benefit is
    customization: the contract can specify the exact metal grade, quantity,
    delivery date, location, and settlement terms. The main risk is that OTC
    contracts create more bilateral counterparty, liquidity, legal, and
    valuation risk than standardized exchange-traded futures.

    \item Separation of front, middle, and back office functions reduces the
    chance that the same people who take risk also measure, approve, settle, and
    report that risk. It creates independent checks on trade entry, valuation,
    limits, collateral, and cash movements.
\end{enumerate}
\end{solution}

\section*{Question 2: Payoff, Profit, and Option Logic \pts{3}}

\begin{solution}
For the call,
\[
    \text{call payoff}=\max(\St-100,0),
    \qquad
    \text{call profit}=\text{call payoff}-7.
\]
For the put,
\[
    \text{put payoff}=\max(100-\St,0),
    \qquad
    \text{put profit}=\text{put payoff}-5.
\]

\begin{center}
\begin{tabular}{@{}ccccc@{}}
\toprule
$\St$ & Call payoff & Call profit & Put payoff & Put profit \\
\midrule
88  & 0  & -7 & 12 & 7  \\
100 & 0  & -7 & 0  & -5 \\
116 & 16 & 9  & 0  & -5 \\
\bottomrule
\end{tabular}
\end{center}

An option can have a positive payoff and still produce a negative profit when
the payoff is smaller than the premium paid. For example, a call payoff of 4
would still imply a profit of $4-7=-3$ for this call buyer.

The buyer of a call has limited downside: the most the buyer can lose is the
premium. The buyer of a forward has no such premium cap on losses; a long
forward loses $K-\St$ when the terminal spot price is below the delivery price,
and that loss can be large.
\end{solution}

\section*{Question 3: Futures Settlement, Margin, and Liquidity Risk \pts{3}}

\begin{solution}
The trader is short
\[
    6\times 5{,}000=30{,}000
\]
bushels. The settlement price increased from 580.25 cents to 592.75 cents, a
change of
\[
    592.75-580.25=12.50
\]
cents per bushel, or USD 0.125 per bushel. A short futures position loses when
the futures price rises, so today's futures profit/loss is
\[
    -30{,}000\times 0.125=-\$3{,}750.
\]

The margin account balance after daily settlement is
\[
    21{,}000-3{,}750=\$17{,}250.
\]
Since USD 17,250 is above the maintenance margin of USD 16,000, there is no
margin call. The required deposit is therefore zero.

Daily settlement can create liquidity risk because futures losses must be paid
as cash immediately, even if the hedge reduces the firm's total economic risk.
The offsetting gain on the underlying exposure may arrive only later, when the
physical purchase or sale occurs.
\end{solution}

\section*{Question 4: Hedge Direction, Basis Risk, and Effective Prices \pts{4}}

\begin{solution}
The firm expects to sell the commodity later. It is hurt by falling spot prices,
so it should use a short hedge: losses on the physical sale price are partly
offset by gains on the short futures position when futures prices fall.

The gain on the short futures position is
\[
    \Fone-\Ftwo=47.80-46.20=1.60.
\]
The effective sale price is the spot sale price plus the futures gain:
\[
    S_2+(\Fone-\Ftwo)=45.90+1.60=47.50.
\]

The closing basis is
\[
    b_2=S_2-\Ftwo=45.90-46.20=-0.30.
\]
Therefore
\[
    \Fone+b_2=47.80-0.30=47.50,
\]
which equals the effective sale price.

If the firm originally expected $b_2=-0.60$, the expected effective sale price
was
\[
    47.80-0.60=47.20.
\]
The realized effective sale price of 47.50 was better by 0.30 per unit. The
basis was stronger than expected.
\end{solution}

\section*{Question 5: Cross Hedging and the Minimum-Variance Hedge Ratio \pts{4}}

\begin{solution}
The minimum-variance hedge ratio is
\[
    h^*=\rho\frac{\sigma_S}{\sigma_F}
    =0.82\frac{0.24}{0.20}
    =0.82(1.20)
    =0.9840.
\]
The unrounded number of contracts is
\[
    N^*=h^*\frac{Q_A}{Q_F}
    =0.9840\frac{1{,}680{,}000}{42{,}000}
    =0.9840(40)
    =39.36.
\]

The nearest integer hedge is 39 contracts. A practical recommendation is to use
39 contracts if the objective is to stay closest to the minimum-variance
quantity. If the airline is especially concerned about under-hedging fuel price
increases, 40 contracts is also defensible, but it is a slight over-hedge
relative to the estimated minimum-variance position.

The airline should go long futures because it will buy jet fuel later and is
hurt by rising fuel prices. A long futures position gains when the futures price
rises.

Residual risk remains because jet fuel and heating oil are different assets, so
their prices may not move one-for-one. Other residual risks include timing
mismatch, location mismatch, changing correlations, basis risk when the hedge is
closed, and contract-size rounding.
\end{solution}

\section*{Question 6: Forward Pricing and Contract Value \pts{3}}

\begin{solution}
For a stock with known cash income, the no-arbitrage forward price is
\[
    \Fzero=(S_0-I)e^{rT}.
\]
Here $S_0=96$, $I=1.80$, $r=0.04$, and $T=0.75$, so
\[
    \Fzero=(96-1.80)e^{0.04\times 0.75}
    =94.20e^{0.03}
    \approx 97.07.
\]

The quoted forward price of 99.50 is too high relative to the no-arbitrage
price. The cash-and-carry arbitrage is:
\begin{enumerate}
    \item Buy the stock today.
    \item Finance the purchase, taking account of the present value of the known
    dividends.
    \item Short the forward at 99.50.
    \item At maturity, deliver the stock into the forward and receive 99.50.
\end{enumerate}
The financing cost net of the known dividends is
\[
    (S_0-I)e^{rT}=97.07.
\]
The locked-in profit per share is therefore approximately
\[
    99.50-97.07=2.43.
\]

The value of an old long forward contract is
\[
    f=(\Fzero-\K)e^{-rT}.
\]
With $\Fzero=94$, $\K=92$, $r=0.04$, and $T=0.5$,
\[
    f=(94-92)e^{-0.04\times 0.5}
    =2e^{-0.02}
    \approx 1.96.
\]
The old long forward is valuable because it allows the holder to buy at 92 when
a new forward with the same maturity would lock in a delivery price of 94.
\end{solution}

\end{document}
